\begin{abstract}
轴承是旋转机械中最常见、也最容易影响整机可靠性的关键部件。为了在失效前完成维护决策，预测性维护不仅需要判断轴承是否已经故障，还需要估计剩余寿命、识别退化阶段，并尽早发现异常趋势。本文围绕 XJTU-SY 和 PHM2012 两个轴承振动数据集，构建了一个覆盖数据加载、特征分析、标签构造、任务数据集、模型训练、对照实验和演示视频的完整实验闭环。

本文将任务统一为三个层次：剩余寿命预测、健康阶段识别和早期故障检测。其中，剩余寿命预测统一采用 \code{linear_rul_norm} 作为目标，HealthState 和 EarlyFault 则基于退化过程构造伪标签。特征侧采用 manual\_basic 特征族，覆盖时域幅值、高阶统计、频域形状、频域复杂度、频带能量和合成幅值通道，并通过相关性、趋势性、阶段可分性和早期故障效应分析筛选任务适配特征。模型侧先完成 MLP、XGBoost、RandomForest 等表格基线，再以长度为 8 的特征序列作为输入训练 GRU 序列基线。最终主线结果来自六个 200ep GRU 序列实验；两个 50ep 加速视频仅用于展示训练过程。

实验结果表明，GRU 序列基线在 XJTU-SY EarlyFault 上取得 test WeightedF1 = 0.851679，在 PHM2012 EarlyFault 上取得 test WeightedF1 = 0.682715；RUL 任务在 XJTU-SY 和 PHM2012 上分别取得 test RMSE = 0.275009 和 0.280362。与此同时，预测质量排查也显示 RUL 泛化仍是主要难点，尤其在跨轴承、跨工况条件下不应过度夸大预测精度。除本文方法实验外，本文还整理了 PHM2012 RUL 的 CAM-LSTM 复现模型和 XJTU-SY 故障诊断复现模型，用于说明典型论文方法的复现实验能力，但不将其混入本文统一三任务主结果。
\end{abstract}

\addcontentsline{toc}{section}{摘要}
\newpage
